\documentclass[11pt,a4paper]{scrartcl}

\usepackage{../../1.\space Vienanariai\space ir\space daugianariai/manostilius_lt}

\title{Nepilnosios kvadratinės lygtys}
\author{Andrius Berniukevičius}

\begin{document}

\maketitle

\section{Kas yra nepilnoji kvadratinė lygtis?}

\begin{definition}
\vocab{Pilnoji kvadratinė lygtis} -- tai lygtis, kurią galima užrašyti pavidalu $ax^2+bx+c=0$, kur $a \neq 0$[cite: 2].
\end{definition}

Jei lygtyje koeficientas $b=0$ arba $c=0$, tokia lygtis vadinama \vocab{nepilnąja kvadratine lygtimi}. Šiame skyriuje nagrinėsime abiejų tipų nepilnąsias lygtis: kai $b=0$ (lygtys pavidalo $ax^2+c=0$) ir kai $c=0$ (lygtys pavidalo $ax^2+bx=0$).

\begin{remark}
Prieš pradedant spręsti lygtis, būtina prisiminti svarbią taisyklę apie kvadratines šaknis. Dažnai daroma klaida manant, kad $\sqrt{9} = \pm 3$. Tai yra netiesa! Pagal kvadratinės šaknies apibrėžimą, aritmetinė kvadratinė šaknis yra \textbf{neneigiamas} skaičius, todėl griežtai $\sqrt{9} = 3$. 

Lygties $x^2 = 9$ sprendiniai yra $x = \pm 3$ dėl to, kad mes ieškome visų skaičių (tiek teigiamų, tiek neigiamų), kuriuos pakėlus kvadratu gauname $9$. Mes tai užrašome kaip $x = \pm\sqrt{9}$, o pati šaknies reikšmė visada išlieka teigiama[cite: 2].
\end{remark}

\section{Lygties $ax^2+c=0$ sprendimas dviem būdais}

Šio tipo lygtis galima spręsti dviem skirtingais, bet lygiaverčiais būdais: suvedant lygtį į pavidalą $x = \pm\sqrt{d}$ arba skaidant reiškinį dauginamaisiais pagal kvadratų skirtumo formulę[cite: 2]. Šiame skyrelyje ieškosime sprendinių \textbf{realiųjų skaičių aibėje} $\mathbb{R}$.

\begin{example}
Išspręskite lygtį realiųjų skaičių aibėje:
\[ x^2 - 25 = 0 \]
\end{example}
\begin{solution}
\textbf{1 būdas: Suvedant į $x = \pm\sqrt{d}$.}\\
Perkeliame laisvąjį narį į dešinę pusę:
\[ x^2 = 25 \]
Ieškome skaičių, kurių kvadratas lygus $25$. Traukiame šaknį, nepamiršdami $\pm$ ženklo prieš ją:
\[ x = \pm\sqrt{25} \]
\[ x = \pm 5 \]

\textbf{2 būdas: Skaidant dauginamaisiais.}\\
Pritaikome kvadratų skirtumo formulę $a^2 - b^2 = (a-b)(a+b)$:
\[ x^2 - 5^2 = 0 \]
\[ (x-5)(x+5) = 0 \]
Sandauga lygi nuliui tada ir tik tada, kai bent vienas iš dauginamųjų lygus nuliui:
\[ x-5 = 0 \quad \text{arba} \quad x+5 = 0 \]
\[ x = 5 \quad \text{arba} \quad x = -5 \]
\textbf{Atsakymas:} $x = \pm 5$.
\end{solution}

\begin{example}
Išspręskite lygtį realiųjų skaičių aibėje:
\[ 2x^2 - 18 = 0 \]
\end{example}
\begin{solution}
\textbf{1 būdas: Suvedant į $x = \pm\sqrt{d}$.}\\
Perkeliame skaičių ir padalijame abi puses iš $2$:
\[
\begin{aligned}
2x^2 &= 18 \\
x^2 &= 9 \\
x &= \pm\sqrt{9} \\
x &= \pm 3.
\end{aligned}
\]

\textbf{2 būdas: Skaidant dauginamaisiais.}\\
Iškeliame bendrą dauginamąjį, o tada pritaikome formulę:
\[
\begin{aligned}
2(x^2 - 9) &= 0 \\
2(x-3)(x+3) &= 0
\end{aligned}
\]
Padaliję iš 2, gauname tuos pačius sprendinius:
\[ x-3 = 0 \quad \text{arba} \quad x+3 = 0 \]
\[ x = 3 \quad \text{arba} \quad x = -3 \]
\textbf{Atsakymas:} $x = \pm 3$.
\end{solution}

\begin{example}
Išspręskite lygtį realiųjų skaičių aibėje:
\[ x^2 + 9 = 0 \]
\end{example}
\begin{solution}
Perkeliame laisvąjį narį į dešinę pusę:
\[ x^2 = -9 \]
Kadangi bet kokio realiojo skaičiaus kvadratas yra neneigiamas ($x^2 \ge 0$), nėra tokio realiojo skaičiaus, kurį pakėlus kvadratu gautume neigiamą rezultatą (šiuo atveju $-9$). 
\textbf{Atsakymas:} realiųjų sprendinių nėra ($\varnothing$).
\end{solution}

\section{Sudėtingesni koeficientai ir iracionalūs sprendiniai ($ax^2+c=0$)}

Kartais lygties koeficientai gali atrodyti sudėtingai, pavyzdžiui, būti išreikšti paprastosiomis arba mišriosiomis trupmenomis. Tačiau atlikus veiksmus, sprendiniai dažnai gaunami labai gražūs (racionalūs skaičiai).

\begin{example}
Išspręskite lygtį realiųjų skaičių aibėje:
\[ 1\frac{1}{3}x^2 - \frac{3}{4} = 0 \]
\end{example}
\begin{solution}
Pirmiausia mišriąją trupmeną paverčiame netaisyklingąja ir perkeliame laisvąjį narį į dešinę pusę:
\[ \frac{4}{3}x^2 = \frac{3}{4} \]
Padalijame abi puses iš $\frac{4}{3}$ (arba padauginame iš $\frac{3}{4}$):
\[
\begin{aligned}
x^2 &= \frac{3}{4} \cdot \frac{3}{4} \\
x^2 &= \frac{9}{16}
\end{aligned}
\]
Traukiame šaknį:
\[ x = \pm\sqrt{\frac{9}{16}} \]
\[ x = \pm\frac{3}{4} \]
\textbf{Atsakymas:} $x = \pm\frac{3}{4}$.
\end{solution}

Kita vertus, ne visada ištraukus šaknį gauname sveiką ar racionalų skaičių. Jei skaičius neturi tikslios šaknies, sprendinį paliekame su šaknies ženklu. Tokie sprendiniai vadinami \textbf{iracionaliaisiais}.

\begin{example}
Išspręskite lygtį realiųjų skaičių aibėje:
\[ 2x^2 - 14 = 0 \]
\end{example}
\begin{solution}
Perkeliame laisvąjį narį ir padalijame iš 2:
\[
\begin{aligned}
2x^2 &= 14 \\
x^2 &= 7
\end{aligned}
\]
Kadangi 7 nėra tikslus kvadratas ir tikslios šaknies neturi, traukdami šaknį gauname iracionaliuosius skaičius:
\[ x = \pm\sqrt{7} \]
\textbf{Atsakymas:} $x = \pm\sqrt{7}$.
\end{solution}

\section{Lygties $ax^2+bx=0$ sprendimas}

Kai lygtyje laisvasis narys $c=0$, gauname lygtį, kurioje abu nariai turi kintamąjį $x$. Tokios lygtys realiųjų skaičių aibėje sprendžiamos \textbf{iškeliant bendrąjį dauginamąjį prieš skliaustus}.

\begin{example}
Išspręskite lygtį realiųjų skaičių aibėje:
\[ x^2 - 7x = 0 \]
\end{example}
\begin{solution}
Iškeliame bendrąjį dauginamąjį $x$ prieš skliaustus:
\[ x(x - 7) = 0 \]
Sandauga lygi nuliui, kai bent vienas iš dauginamųjų lygus nuliui:
\[ x = 0 \quad \text{arba} \quad x - 7 = 0 \]
Iš antrosios lygties gauname:
\[ x = 7 \]
\textbf{Atsakymas:} $x = 0$; $x = 7$.
\end{solution}

Jei lygties koeficientai yra trupmeniniai, sprendimo eiga išlieka lygiai tokia pati – iškeliame kintamąjį (ir, jei patogu, trupmeninį koeficientą) prieš skliaustus.

\begin{example}
Išspręskite lygtį realiųjų skaičių aibėje:
\[ \frac{2}{3}x^2 - 2,4x = 0 \]
\end{example}
\begin{solution}
Pirmiausia dešimtainį skaičių $2,4$ užrašome paprastosios trupmenos pavidalu:
\[ 2,4 = \frac{12}{5} \]
Taigi lygtis tampa:
\[ \frac{2}{3}x^2 - \frac{12}{5}x = 0 \]
Iškeliame bendrąjį dauginamąjį $x$ prieš skliaustus:
\[ x\left( \frac{2}{3}x - \frac{12}{5} \right) = 0 \]
Sandauga lygi nuliui, kai bent vienas dauginamasis lygus nuliui:
\[ x = 0 \quad \text{arba} \quad \frac{2}{3}x - \frac{12}{5} = 0 \]
Sprendžiame antrąją lygtį:
\[ \frac{2}{3}x = \frac{12}{5} \]
\[ x = \frac{12}{5} \cdot \frac{3}{2} = \frac{36}{10} = 3,6 \]
\textbf{Atsakymas:} $x = 0$; $x = 3,6$.
\end{solution}

\section{Kompleksiniai sprendiniai: lygties $x^2 = -a$ atvejis}

Kas nutinka, jei sprendžiant lygtį $ax^2+c=0$, dešinėje pusėje gauname neigiamą skaičių? Kaip ką tik matėme pavyzdyje su lygtimi $x^2 + 9 = 0$, \textbf{realiųjų skaičių aibėje} nerasime skaičiaus, kurio kvadratas būtų neigiamas, todėl tokia lygtis sprendinių neturi. Kad išspręstume tokias lygtis, mums reikia \textbf{kompleksinių skaičių aibės} $\mathbb{C}$.

Naudojant kompleksinius skaičius, lygtis, kurios anksčiau neturėjo sprendinių realiųjų skaičių aibėje, tampa lengvai išsprendžiamos.

\begin{example}
Išspręskite prieš tai nagrinėtą lygtį kompleksinių skaičių aibėje:
\[ x^2 + 9 = 0 \]
\end{example}
\begin{solution}
Perkeliame $9$ į dešinę pusę:
\[ x^2 = -9 \]
Kaip jau matėme, realiųjų skaičių aibėje sprendinių nėra. Tačiau kompleksinių skaičių aibėje vietoje minuso ženklo įvedame $i^2$:
\[
\begin{aligned}
x^2 &= 9 \cdot (-1) \\
x^2 &= 9i^2
\end{aligned}
\]
Dabar traukiame kvadratinę šaknį, vėl naudodami $\pm$:
\[ x = \pm\sqrt{9i^2} \]
Kadangi $\sqrt{9} = 3$ ir $\sqrt{i^2} = i$, gauname:
\[ x = \pm 3i \]
\textbf{Atsakymas:} $x = \pm 3i$.
\end{solution}

\begin{example}
Išspręskite lygtį kompleksinių skaičių aibėje:
\[ 2x^2 + 50 = 0 \]
\end{example}
\begin{solution}
Pertvarkome lygtį:
\[
\begin{aligned}
2x^2 &= -50 \\
x^2 &= -25 \\
x^2 &= 25i^2 \\
x &= \pm\sqrt{25i^2} \\
x &= \pm 5i.
\end{aligned}
\]
\textbf{Atsakymas:} $x = \pm 5i$.
\end{solution}

\section{Lygties $(ax+b)^2=c$ sprendimas dviem būdais}

Lygtį, kurios kairėje pusėje yra dvinario kvadratas, galima spręsti dviem būdais. Pirmuoju būdu ištraukiame kvadratinę šaknį ir prisimename, kad reikia užrašyti abu ženklus. Antruoju būdu išskleidžiame kvadratą, perkeliame visus narius į vieną pusę ir pritaikome kvadratų skirtumo formulę.

\begin{example}
Išspręskite lygtį realiųjų skaičių aibėje dviem būdais:
\[ (2x-3)^2=25. \]
\end{example}
\begin{solution}
\textbf{1 būdas: Ištraukiant kvadratinę šaknį.}

Ištraukiame kvadratinę šaknį iš abiejų lygties pusių:
\[
2x-3=\pm\sqrt{25}=\pm5.
\]
Todėl gauname dvi lygtis:
\[
\begin{aligned}
2x-3&=5, & 2x&=8, & x&=4,\\
2x-3&=-5, & 2x&=-2, & x&=-1.
\end{aligned}
\]

\textbf{2 būdas: Skaidant dauginamaisiais.}

Išskleidžiame kvadratą ir perkeliame visus narius į kairę pusę:
\[
\begin{aligned}
(2x-3)^2-25&=0,\\
\bigl((2x-3)-5\bigr)\bigl((2x-3)+5\bigr)&=0,\\
(2x-8)(2x+2)&=0.
\end{aligned}
\]
Taigi
\[
2x-8=0 \quad \text{arba} \quad 2x+2=0,
\]
todėl $x=4$ arba $x=-1$.

\textbf{Atsakymas:} $x=4$ arba $x=-1$.
\end{solution}

\begin{example}
Išspręskite lygtį realiųjų skaičių aibėje dviem būdais:
\[ (3x+1)^2=16. \]
\end{example}
\begin{solution}
\textbf{1 būdas: Ištraukiant kvadratinę šaknį.}

\[
3x+1=\pm\sqrt{16}=\pm4.
\]
Todėl
\[
\begin{aligned}
3x+1&=4, & 3x&=3, & x&=1,\\
3x+1&=-4, & 3x&=-5, & x&=-\frac{5}{3}.
\end{aligned}
\]

\textbf{2 būdas: Skaidant dauginamaisiais.}

\[
\begin{aligned}
(3x+1)^2-16&=0,\\
\bigl((3x+1)-4\bigr)\bigl((3x+1)+4\bigr)&=0,\\
(3x-3)(3x+5)&=0.
\end{aligned}
\]
Vadinasi,
\[
3x-3=0 \quad \text{arba} \quad 3x+5=0,
\]
todėl $x=1$ arba $x=-\frac{5}{3}$.

\textbf{Atsakymas:} $x=1$ arba $x=-\frac{5}{3}$.
\end{solution}

\section{Uždaviniai}

\begin{problem}
Išspręskite nepilnąsias kvadratines lygtis, skaidydami dauginamaisiais pagal kvadratų skirtumo formulę, realiųjų skaičių aibėje:
\begin{tasks}[label={(\arabic*)}, label-offset=0.9em](2)
    \task $x^2-16=0$
    \task $x^2-64=0$
    \task $x^2 + 9 = 0$
    \task $5x^2-20=0$
    \task $x^2 - 5 = 0$
    \task $-x^2 + 81 = 0$
    \task $-3x^2 + 27 = 0$
    \task $3x^2 - 33 = 0$
    \task $5x^2 - 15 = 0$
    \task $2x^2 - 10 = 0$
    \task $3x^2-27=0$
    \task $2x^2 + 7 = 0$
\end{tasks}
\end{problem}

\begin{problem}
Išspręskite tas pačias nepilnąsias kvadratines lygtis, suvedant jas į pavidalą $x = \pm\sqrt{d}$, realiųjų skaičių aibėje:
\begin{tasks}[label={(\arabic*)}, label-offset=0.9em](2)
    \task $x^2-16=0$
    \task $x^2-64=0$
    \task $x^2 + 9 = 0$
    \task $5x^2-20=0$
    \task $x^2 - 5 = 0$
    \task $-x^2 + 81 = 0$
    \task $-3x^2 + 27 = 0$
    \task $3x^2 - 33 = 0$
    \task $5x^2 - 15 = 0$
    \task $2x^2 - 10 = 0$
    \task $3x^2-27=0$
    \task $2x^2 + 7 = 0$
\end{tasks}
\end{problem}

\begin{problem}
Išspręskite nepilnąsias kvadratines lygtis su trupmeniniais koeficientais realiųjų skaičių aibėje:
\begin{tasks}[label={(\arabic*)}, label-offset=0.9em](2)
    \task $\frac{1}{2}x^2 - 8 = 0$
    \task $\frac{2}{5}x^2 - 3\frac{3}{5} = 0$
    \task $1\frac{7}{9}x^2 - \frac{4}{25} = 0$
    \task $\frac{3}{7}x^2 - \frac{27}{28} = 0$
\end{tasks}
\end{problem}

\begin{problem}
Išspręskite lygtis realiųjų skaičių aibėje dviem būdais:
\begin{tasks}[label={(\arabic*)}, label-offset=0.9em](2)
    \task $(2x+1)^2=9$
    \task $(3x-2)^2=49$
    \task $(x+4)^2=25$
    \task $(2x-5)^2=81$
    \task $(5x+2)^2=16$
    \task $(4x-3)^2=36$
\end{tasks}
\end{problem}

\begin{problem}
Išspręskite lygtis kompleksinių skaičių aibėje $\mathbb{C}$:
\begin{tasks}[label={(\arabic*)}, label-offset=0.9em](4)
    \task $x^2 = -4$
    \task $x^2 = -36$
    \task $x^2+81=0$
    \task $x^2+100=0$
    \task $2x^2+18=0$
    \task $4x^2+64=0$
    \task $10x^2 = -1000$
    \task $\frac{1}{2}x^2+8=0$
\end{tasks}
\end{problem}

\newpage 
\begin{problem}
Išspręskite nepilnąsias kvadratines lygtis, iškeldami bendrąjį dauginamąjį prieš skliaustus, realiųjų skaičių aibėje:
\begin{tasks}[label={(\arabic*)}, label-offset=0.9em](2)
    \task $x^2-5x=0$
    \task $x^2+8x=0$
    \task $2x^2-10x=0$
    \task $7x^2+14x=0$
    \task $\frac{1}{3}x^2 - \frac{4}{3}x = 0$
    \task $\frac{3}{4}x^2 + 1\frac{1}{2}x = 0$
    \task $1\frac{1}{5}x^2 - \frac{3}{5}x = 0$
    \task $\frac{2}{7}x^2 + \frac{4}{7}x = 0$
\end{tasks}
\end{problem}

\newpage
\answerssection

\begin{problem} \phantom{text}
\begin{tasks}[label={(\arabic*)}, label-offset=0.9em](2)
    \task $x=\pm 4$;
    \task $x=\pm 8$;
    \task Realių sprendinių nėra;
    \task $x=\pm 2$;
    \task $x=\pm \sqrt{5}$;
    \task $x=\pm 9$;
    \task $x=\pm 3$;
    \task $x=\pm \sqrt{11}$;
    \task $x=\pm \sqrt{3}$;
    \task $x=\pm \sqrt{5}$;
    \task $x=\pm 3$;
    \task Realių sprendinių nėra.
\end{tasks}
\end{problem}

\begin{problem} \phantom{text}
\begin{tasks}[label={(\arabic*)}, label-offset=0.9em](2)
    \task $x=\pm 4$;
    \task $x=\pm 8$;
    \task Realių sprendinių nėra;
    \task $x=\pm 2$;
    \task $x=\pm \sqrt{5}$;
    \task $x=\pm 9$;
    \task $x=\pm 3$;
    \task $x=\pm \sqrt{11}$;
    \task $x=\pm \sqrt{3}$;
    \task $x=\pm \sqrt{5}$;
    \task $x=\pm 3$;
    \task Realių sprendinių nėra.
\end{tasks}
\end{problem}

\begin{problem} \phantom{text}
\begin{tasks}[label={(\arabic*)}, label-offset=0.9em](2)
    \task $x=\pm 4$;
    \task $x=\pm 3$;
    \task $x=\pm \frac{3}{10}$;
    \task $x=\pm \frac{3}{2}$.
\end{tasks}
\end{problem}

\begin{problem} \phantom{text}
\begin{tasks}[label={(\arabic*)}, label-offset=0.9em](2)
    \task $x=1$, $x=-2$;
    \task $x=3$, $x=-\frac{5}{3}$;
    \task $x=1$, $x=-9$;
    \task $x=7$, $x=-2$;
    \task $x=\frac{2}{5}$, $x=-\frac{6}{5}$;
    \task $x=\frac{9}{4}$, $x=-\frac{3}{4}$.
\end{tasks}
\end{problem}

\begin{problem} \phantom{text}
\begin{tasks}[label={(\arabic*)}, label-offset=0.9em](4)
    \task $x=\pm 2i$;
    \task $x=\pm 6i$;
    \task $x=\pm 9i$;
    \task $x=\pm 10i$;
    \task $x=\pm 3i$;
    \task $x=\pm 4i$;
    \task $x=\pm 10i$;
    \task $x=\pm 4i$.
\end{tasks}
\end{problem}

\begin{problem} \phantom{text}
\begin{tasks}[label={(\arabic*)}, label-offset=0.9em](4)
    \task $x=0$, $x=5$;
    \task $x=0$, $x=-8$;
    \task $x=0$, $x=5$;
    \task $x=0$, $x=-2$;
    \task $x=0$, $x=4$;
    \task $x=0$, $x=-2$;
    \task $x=0$, $x=\frac{1}{2}$;
    \task $x=0$, $x=-2$.
\end{tasks}
\end{problem}

\end{document}